\section{PHÂN TÍCH CÁC YẾU TỐ TÁC ĐỘNG VÀ ĐÚC KẾT BÀI HỌC KINH NGHIỆM LỊCH SỬ}

Nhìn nhận lịch sử dưới lăng kính khoa học giáo dục và thể chế học đòi hỏi nghiên cứu sinh không chỉ dừng lại ở việc phục dựng các sự kiện niên biểu, mà phải đi sâu mổ xẻ các mối quan hệ nhân quả, bóc tách các nhóm yếu tố khách quan và chủ quan chi phối sự vận động của hệ thống, đồng thời đúc kết những bài học kinh nghiệm mang tính quy luật có thể soi đường cho hiện tại và tương lai. Chương này tập trung phân tích các yếu tố tác động, xây dựng ma trận đối sánh toàn diện sáu giai đoạn lịch sử và luận giải bốn bài học cốt lõi rút ra từ chặng đường hơn tám mươi năm qua.

\subsection{Các nhóm yếu tố kinh tế, xã hội và chính trị tác động đến giáo dục nghề nghiệp}

Sự hình thành, phát triển hay thoái trào của giáo dục nghề nghiệp luôn là tấm gương phản chiếu chân thực cấu trúc kinh tế và trật tự xã hội của quốc gia. Qua rà soát tài liệu lịch sử, có thể khái quát ba nhóm yếu tố căn bản đã tác động trực diện đến hệ sinh thái đào tạo nghề Việt Nam qua các thời kỳ:

\subsubsection{Yếu tố chuyển đổi mô hình kinh tế và chuyển dịch cơ cấu việc làm}

Giáo dục nghề nghiệp luôn là bộ phận nhạy cảm nhất đối với những thay đổi của mô hình kinh tế. Trong thời kỳ kháng chiến cứu nước (1945 -- 1975), nền kinh tế phục vụ chiến tranh đặt mục tiêu tự cấp tự túc vũ khí lên hàng đầu, do đó đào tạo nghề mang tính chất cơ động, ưu tiên tuyệt đối cho cơ khí quân giới và vận tải. Bước sang thời kỳ kinh tế kế hoạch hóa tập trung (1975 -- 1986), sự tồn tại của khu vực kinh tế quốc doanh bao cấp đã tạo ra một thị trường lao động định sẵn, trong đó nhà nước đóng vai trò vừa là người đặt hàng đào tạo duy nhất, vừa là cơ quan sử dụng lao động độc quyền.

Khi đất nước chuyển đổi sang nền kinh tế thị trường định hướng xã hội chủ nghĩa từ sau Đại hội Đảng lần thứ VI năm 1986, các quy luật giá trị, quy luật cung cầu và quy luật cạnh tranh đã lập tức làm thay đổi hoàn toàn sân chơi của các trường nghề. Doanh nghiệp không còn tiếp nhận lao động theo chỉ tiêu hành chính mà đòi hỏi tay nghề thực chất đáp ứng ngay vị trí việc làm. 

Đặc biệt, sự dịch chuyển mạnh mẽ của cơ cấu kinh tế từ thuần túy nông nghiệp sang công nghiệp chế biến, chế tạo, xây dựng, dịch vụ và logistics trong ba thập kỷ qua đã tạo ra sự biến động to lớn về cơ cấu việc làm. Sự xuất hiện của các tập đoàn đa quốc gia và dòng vốn đầu tư trực tiếp nước ngoài đòi hỏi người lao động kỹ thuật không chỉ có kỹ năng thao tác tay chân đơn thuần mà phải làm chủ được công nghệ tự động hóa, kỹ năng ngoại ngữ và tác phong công nghiệp hiện đại.

\subsubsection{Yếu tố tâm lý xã hội và định kiến bằng cấp trong văn hóa Á Đông}

Một trong những rào cản văn hóa -- xã hội nặng nề và dai dẳng nhất kìm hãm sự phát triển của giáo dục nghề nghiệp Việt Nam chính là tâm lý sùng bái bằng cấp hàn lâm, bắt nguồn sâu xa từ truyền thống khoa cử phong kiến kéo dài hàng ngàn năm. Trong tâm thức của phần lớn các bậc phụ huynh và học sinh Việt Nam, con đường đại học được coi là con đường duy nhất để tiến thân, lập nghiệp và khẳng định địa vị xã hội; trong khi đó, việc vào học tại các trường dạy nghề, trung cấp hay cao đẳng nghề thường bị xem là sự lựa chọn thứ cấp, là giải pháp tình thế bất đắc dĩ khi học sinh thi trượt đại học.

Định kiến xã hội này đã dẫn đến hiện tượng tâm lý lệch lạc kéo dài qua nhiều thế hệ: người ta sẵn sàng chi tiêu tốn kém để con em theo học các ngành đại học xã hội, kinh tế chung chung dù cơ hội việc làm rất bấp bênh, hơn là cho con đi học nghề kỹ thuật có nhu cầu tuyển dụng thực tế rất cao và thu nhập ổn định. Hệ quả là công tác phân luồng học sinh sau trung học cơ sở và trung học phổ thông luôn là bài toán nan giải đối với các nhà quản lý, dẫn đến sự mất cân đối nghiêm trọng trong cơ cấu nguồn nhân lực quốc gia.

\subsubsection{Yếu tố thể chế, chính sách và nguồn lực đầu tư công}

Sự phát triển của giáo dục nghề nghiệp phụ thuộc mang tính quyết định vào sự nhất quán của thể chế và mức độ ưu tiên của chính sách công. Lịch sử cho thấy những giai đoạn nào bộ máy quản lý nhà nước được tổ chức độc lập, tập trung và có thẩm quyền mạnh mẽ (như thời kỳ Tổng cục Đào tạo công nhân kỹ thuật trực thuộc Bộ Lao động từ năm 1969, hay Tổng cục Dạy nghề trực thuộc Hội đồng Chính phủ giai đoạn 1978 -- 1987) thì vị thế của ngành nghề được tôn vinh, mạng lưới trường lớp được đầu tư bài bản và thu hút được nhiều nhân tài. 

Ngược lại, trong những giai đoạn chức năng quản lý bị phân tán, hạ cấp xuống cấp vụ hoặc bị chia cắt song trùng giữa hai bộ quản lý (giai đoạn 2006 -- 2014), hệ thống rơi vào tình trạng lúng túng, nguồn lực đầu tư bị phân tán và chính sách thiếu tính đồng bộ. 

Mặt khác, mức độ đầu tư công cho cơ sở vật chất, nhà xưởng thực hành và chế độ đãi ngộ cho nhà giáo giáo dục nghề nghiệp luôn là nhân tố quyết định chất lượng đầu ra. Tại các cơ sở đào tạo nghề kỹ thuật, chi phí đào tạo trên mỗi học sinh luôn cao gấp nhiều lần so với đào tạo các ngành khoa học xã hội do đòi hỏi máy móc, nguyên vật liệu thực hành tiêu hao đắt tiền. Sự thiếu hụt kinh phí công và cơ chế tài chính xơ cứng trong nhiều thời kỳ đã khiến trang thiết bị của nhiều trường nghề chậm đổi mới, tụt hậu xa so với dây chuyền sản xuất của doanh nghiệp ngoài thị trường.

\subsection{Đánh giá tổng hợp thành tựu và hạn chế lịch sử qua bảng ma trận so sánh}

\subsubsection{Những thành tựu có ý nghĩa chiến lược của hệ thống}

Trải qua hơn tám thập kỷ đồng hành cùng sự nghiệp cách mạng, hệ thống giáo dục nghề nghiệp Việt Nam đã đạt được những thành tựu to lớn, để lại dấu ấn sâu đậm trong lịch sử dân tộc:

Thành tựu lớn nhất là đã đào tạo và cung ứng cho đất nước hàng chục triệu lượt công nhân kỹ thuật, người lao động có kỹ năng qua các thời kỳ. Lực lượng này chính là nòng cốt xây dựng nên các công trình công nghiệp biểu tượng của đất nước như Nhà máy Gang thép Thái Nguyên, Thủy điện Thác Bà, Thủy điện Hòa Bình, cầu Thăng Long, các giàn khoan dầu khí ngoài khơi Vũng Tàu, và hiện nay là lực lượng vận hành các tổ hợp công nghiệp công nghệ cao của Samsung, Foxconn, VinFast.

Thành tựu thứ hai là đã xây dựng được một mạng lưới cơ sở giáo dục nghề nghiệp rộng khắp trên phạm vi toàn quốc, từ các trung tâm giáo dục nghề nghiệp cấp huyện đến các trường cao đẳng chất lượng cao cấp tỉnh và cấp vùng. Mạng lưới này không chỉ đào tạo thanh niên học nghề chính quy mà còn thực hiện chức năng an sinh xã hội sâu sắc, đào tạo nghề miễn phí cho hàng triệu lao động nông thôn, con em đồng bào dân tộc thiểu số, bộ đội xuất ngũ và người khuyết tật, góp phần quan trọng vào công cuộc xóa đói giảm nghèo và chuyển dịch cơ cấu kinh tế nông thôn.

Thành tựu thứ ba là từng bước chuẩn hóa chất lượng đào tạo và hội nhập khu vực. Đội ngũ học sinh, sinh viên Việt Nam tham gia các Kỳ thi kỹ năng nghề ASEAN và Kỳ thi kỹ năng nghề thế giới đã liên tục đạt được nhiều huy chương vàng, bạc danh giá trong các nghề cơ điện tử, thiết kế đồ họa, tự động hóa công nghiệp và hàn công nghệ cao, khẳng định tiềm năng và bàn tay vàng của người thợ Việt Nam hoàn toàn có thể vươn tầm đỉnh cao thế giới.

\subsubsection{Các điểm nghẽn cố hữu và hạn chế mang tính hệ thống}

Bên cạnh những thành tựu vẻ vang, rà soát lịch sử cũng chỉ rõ những hạn chế mang tính hệ thống và các điểm nghẽn kéo dài chưa được tháo gỡ triệt để:

Thứ nhất là sự mất cân đối nghiêm trọng trong cơ cấu trình độ nhân lực quốc gia. Trong khi tỷ lệ chuẩn của các nền kinh tế công nghiệp phát triển thường có hình tháp đáy rộng với tỷ lệ một đại học -- bốn cao đẳng, trung cấp -- mười công nhân lành nghề, thì tại Việt Nam, cơ cấu nhân lực đào tạo qua nhiều thập kỷ bị biến dạng theo hướng "tháp lộn ngược" hoặc "tháp phình đầu". Số lượng người học đại học luôn vượt trội so với số người học cao đẳng và trung cấp nghề, dẫn đến tình trạng vừa thừa cử nhân đại học thiếu việc làm, vừa khan hiếm trầm trọng thợ kỹ thuật lành nghề tại các khu công nghiệp.

Thứ hai là sự gắn kết giữa cơ sở giáo dục nghề nghiệp với doanh nghiệp còn mang tính hình thức, thiếu cơ chế ràng buộc trách nhiệm pháp lý chặt chẽ. Doanh nghiệp phần lớn vẫn đứng ngoài quá trình xây dựng chương trình đào tạo, chưa chủ động đầu tư tài chính cho nhà trường mà chủ yếu chỉ tiếp nhận học sinh thực tập vào những tháng cuối khóa để tận dụng lao động giá rẻ. Mô hình đào tạo song hành (đào tạo kép) kiểu Đức dù được đưa vào thí điểm từ lâu nhưng chưa thể nhân rộng do thiếu hành lang ưu đãi thuế thực chất cho doanh nghiệp.

Thứ ba là công tác dự báo nhu cầu thị trường lao động còn yếu kém và lạc hậu. Việc mở ngành nghề đào tạo tại nhiều trường vẫn chạy theo năng lực sẵn có của giáo viên hơn là bám sát tín hiệu thị trường, dẫn đến tình trạng nhiều nghề đào tạo ra không có người tuyển dụng, trong khi các ngành kỹ thuật mới như công nghệ bán dẫn, tự động hóa, logistics, năng lượng tái tạo lại thiếu hụt nguồn tuyển sinh nghiêm trọng.

\subsubsection{Bảng ma trận so sánh toàn diện sáu giai đoạn phát triển lịch sử}

Để khái quát bức tranh toàn cảnh và cung cấp cái nhìn so sánh đối sánh trực quan giữa các thời kỳ theo đúng yêu cầu học thuật của học phần, Bảng 1 tổng hợp năm chiều kích căn bản của hệ thống giáo dục nghề nghiệp Việt Nam từ năm 1945 đến nay.

\begin{small}
\setlength{\tabcolsep}{4.5pt}
\begin{longtable}{|p{2.0cm}|p{2.5cm}|p{2.4cm}|p{3.1cm}|p{2.7cm}|}
\caption{\textbf{Ma trận so sánh đặc trưng hệ thống giáo dục nghề nghiệp Việt Nam qua sáu giai đoạn lịch sử}} \label{tab:so_sanh_lich_su} \\
\hline
\textbf{Giai đoạn} & \textbf{Cơ quan chủ quản} & \textbf{Cấp bậc đào tạo} & \textbf{Cơ chế tài chính \& Việc làm} & \textbf{Hạn chế cốt lõi} \\
\hline
\endfirsthead
\multicolumn{5}{c}%
{{\bfseries \tablename\ \thetable{} -- Tiếp theo từ trang trước}} \\
\hline
\textbf{Giai đoạn} & \textbf{Cơ quan chủ quản} & \textbf{Cấp bậc đào tạo} & \textbf{Cơ chế tài chính \& Việc làm} & \textbf{Hạn chế cốt lõi} \\
\hline
\endhead
\hline \multicolumn{5}{|r|}{{Còn tiếp sang trang sau}} \\ \hline
\endfoot
\hline
\endlastfoot

\textbf{1945 -- 1954} & Nha Kỹ thuật vụ (Bộ Quốc gia Giáo dục) & Lớp tập huấn kỹ thuật, trường thực hành nghề & Nhà nước đài thọ tối thiểu; tự sản xuất tự túc trong chiến khu; phục vụ quân giới & Quy mô nhỏ, phân tán, thiếu thốn máy móc, giáo trình sơ khai \\
\hline
\textbf{1954 -- 1975} & Vụ Quản lý nhân công $\rightarrow$ Tổng cục Đào tạo công nhân kỹ thuật & Công nhân kỹ thuật bậc 2/7, 3/7; sơ cấp, trung cấp kỹ thuật & Bao cấp toàn diện tiền ăn, ở, học bổng; phân công công tác về xí nghiệp quốc doanh & Chia cắt hai miền; chỉ tập trung phát triển tại miền Bắc; chiến tranh tàn phá \\
\hline
\textbf{1975 -- 1986} & Tổng cục Dạy nghề (Hội đồng Chính phủ) & Công nhân kỹ thuật các ngành công nghiệp, xây dựng, nông nghiệp & Bao cấp tuyệt đối theo chỉ tiêu pháp lệnh; phân bổ việc làm một trăm phần trăm & Cơ chế xơ cứng, đào tạo xa rời thực tiễn xí nghiệp, học sinh ỷ lại bao cấp \\
\hline
\textbf{1986 -- 2006} & Bộ Đại học, Trung học chuyên nghiệp và Dạy nghề $\rightarrow$ Bộ Giáo dục và Đào tạo $\rightarrow$ Tổng cục Dạy nghề & Trường dạy nghề, trung tâm dạy nghề, trường trung học chuyên nghiệp & Chấm dứt bao cấp việc làm; chuyển sang xã hội hóa, thu học phí, tự chủ một phần & Khủng hoảng tuyển sinh đầu đổi mới; luân chuyển cơ quan quản lý gây xáo trộn \\
\hline
\textbf{2006 -- 2014} & Tổng cục Dạy nghề (Bộ Lao động -- Thương binh và Xã hội) & Ba cấp trình độ: Sơ cấp nghề, Trung cấp nghề, Cao đẳng nghề & Ngân sách hỗ trợ có mục tiêu; các trường thu học phí; triển khai Đề án 1956 nông thôn & Chia cắt song trùng giữa dạy nghề và trung cấp chuyên nghiệp, lãng phí nguồn lực \\
\hline
\textbf{2014 -- nay} & Tổng cục Giáo dục nghề nghiệp (Bộ Lao động -- Thương binh và Xã hội) & Ba trình độ: Sơ cấp, Trung cấp, Cao đẳng (theo Khung trình độ quốc gia) & Đẩy mạnh tự chủ tài chính; liên kết doanh nghiệp; đặt hàng đào tạo chất lượng cao & Phân luồng sau trung học cơ sở còn thấp; tâm lý chuộng đại học; tự chủ tài chính còn vướng \\
\hline
\end{longtable}
\end{small}

\subsection{Bốn bài học kinh nghiệm lịch sử có giá trị quy luật}

Tổng kết chặng đường lịch sử từ năm 1945 đến nay cho phép đúc kết bốn bài học kinh nghiệm mang tính quy luật, có ý nghĩa phương pháp luận sâu sắc đối với công cuộc hoàn thiện hệ thống giáo dục nghề nghiệp hiện nay:

\subsubsection{Bài học thứ nhất: Tính thống nhất, đồng bộ và ổn định của thể chế quản lý nhà nước}

Lịch sử biến chuyển của cơ quan quản lý nhà nước về đào tạo nghề (từ Nha Kỹ thuật vụ năm 1946, qua Tổng cục Đào tạo công nhân kỹ thuật năm 1969, Tổng cục Dạy nghề trực thuộc Chính phủ năm 1978, đến việc sáp nhập vào Bộ Giáo dục và Đào tạo những năm 1990, rồi lại tái lập Tổng cục Dạy nghề thuộc Bộ Lao động -- Thương binh và Xã hội năm 1998, và sau đó hợp nhất theo Luật Giáo dục nghề nghiệp năm 2014) cho thấy một chân lý thực tiễn: \textbf{sự phân tán thể chế là kẻ thù của hiệu quả phát triển}. 

Mỗi lần cơ quan quản lý bị chia cắt, luân chuyển hoặc hạ cấp, hệ thống trường nghề đều phải trải qua những năm tháng xáo trộn tổ chức, đứt gãy về chính sách và suy giảm nguồn lực đầu tư. Bài học kinh nghiệm cốt tử đặt ra là Nhà nước phải duy trì một cơ chế quản lý nhà nước thống nhất, ổn định lâu dài, có đủ thẩm quyền điều phối liên ngành giữa giáo dục, lao động, tài chính và công nghiệp, tránh tình trạng "trống đánh xuôi kèn thổi ngược" giữa các bộ, ngành.

\subsubsection{Bài học thứ hai: Tính cân đối trong quy hoạch cơ cấu nhân lực quốc gia}

Quy luật phát triển kinh tế chỉ ra rằng một nền sản xuất hùng mạnh đòi hỏi một cơ cấu nhân lực hài hòa, trong đó lực lượng lao động trực tiếp có kỹ năng thực hành phải chiếm tỷ trọng áp đảo. Bài học lịch sử của giai đoạn 1954 -- 1975 cho thấy việc Nhà nước quy hoạch chặt chẽ chỉ tiêu đào tạo công nhân kỹ thuật tương xứng với chỉ tiêu kỹ sư đại học đã tạo nên bộ khung vững chắc cho công cuộc công nghiệp hóa miền Bắc. 

Ngược lại, sự bùng nổ mất kiểm soát của các trường đại học trong hai thập kỷ gần đây, cùng với việc hạ chuẩn tuyển sinh đại học quá mức đã hút cạn nguồn tuyển sinh của các trường nghề, tạo ra một tháp nhân lực dị tật thừa thầy thiếu thợ. Bài học rút ra là công tác quy hoạch mạng lưới giáo dục quốc dân phải được điều tiết bằng các công cụ vĩ mô mạnh mẽ của Nhà nước, kiên quyết định hướng quy mô tuyển sinh đại học phù hợp với nhu cầu nghiên cứu, đồng thời tập trung nguồn lực phát triển giáo dục nghề nghiệp để đảm bảo năng lực sản xuất thực tế của nền kinh tế.

\subsubsection{Bài học thứ ba: Sự gắn kết sống còn giữa cơ sở đào tạo nghề với doanh nghiệp}

Giáo dục nghề nghiệp không thể tồn tại và phát triển nếu tách rời khỏi môi trường sản xuất thực tế. Trong thời kỳ bao cấp, sự gắn kết này diễn ra một cách tự nhiên thông qua cơ chế xí nghiệp mẹ đỡ đầu cho trường con (các trường công nhân kỹ thuật trực thuộc trực tiếp các tổng công ty, nhà máy). Khi bước sang kinh tế thị trường, sợi dây liên kết hành chính bị cắt đứt, nếu nhà trường không chủ động tìm đến doanh nghiệp và ngược lại nếu Nhà nước không có chính sách khuyến khích doanh nghiệp tham gia đào tạo, thì giáo dục nghề nghiệp sẽ rơi vào ngõ cụt đào tạo ra những sản phẩm lỗi thời mà thị trường không chấp nhận. 

Bài học lịch sử khẳng định rằng doanh nghiệp phải được xem là một chủ thể đào tạo bình đẳng bên cạnh nhà trường. Cần phải thể chế hóa trách nhiệm xã hội và quyền lợi kinh tế của doanh nghiệp thông qua các chính sách miễn giảm thuế, trừ chi phí đào tạo vào chi phí hợp lý của doanh nghiệp khi họ đầu tư trang thiết bị hoặc tiếp nhận học sinh đến thực tập.

\subsubsection{Bài học thứ tư: Phân luồng học sinh phổ thông sớm và chính sách an sinh xã hội}

Kinh nghiệm triển khai Đề án 1956 về đào tạo nghề cho lao động nông thôn và các chính sách miễn giảm học phí cho học sinh học trung cấp sau khi tốt nghiệp trung học cơ sở (theo Nghị định số 81/2021/NĐ-CP của Chính phủ) chứng minh rằng: \textbf{chính sách tài chính và an sinh xã hội là đòn bẩy hiệu quả nhất để giải quyết bài toán phân luồng}. 

Để thay đổi tâm lý chuộng bằng cấp của người dân, chỉ tuyên truyền vận động đơn thuần là chưa đủ, mà phải có những cú hích kinh tế thiết thực: người học nghề được miễn học phí, được hỗ trợ vay vốn lập nghiệp, thời gian học ngắn, ra trường có việc làm ngay với thu nhập thỏa đáng và có cơ hội liên thông lên các trình độ cao hơn. Bài học lịch sử chỉ ra rằng công tác hướng nghiệp và phân luồng phải được tiến hành bài bản ngay từ bậc trung học cơ sở, giúp học sinh nhận thức đúng đắn về năng lực bản thân và giá trị cao quý của lao động nghề nghiệp.
